% !PRINTED 495
\weberpage{495}

%<<CONT
as the simplest elements of computation. \gesperrt{Kronecker}\textsuperscript{1)} has struck out on a
path different from these.

The theory of \gesperrt{Dedekind} has been presented by its founder
comprehensively, and with ever-increasing simplicity and perfection, in the
last supplement of the three most recent editions of \gesperrt{Dirichlet's}
Vorlesungen \"uber Zahlentheorie.

The theory of \gesperrt{Kronecker} first became known to the wider circle
of mathematicians in the year 1882, through the Festschrift for
\gesperrt{Kummer's} jubilee, and even now it is still not quite easily
accessible\textsuperscript{2)}.

As regards the \gesperrt{Dedekind} theory, then, we may refer the reader
who wishes to penetrate more deeply into this interesting subject to the
works named. Here, on the other hand, the attempt shall be made to
establish the connection between the various theories, and, on the basis
of the \gesperrt{Kronecker} assumptions and means, to develop from them as
much as is necessary for the algebraic applications that follow.

\begin{center}
\S.\,135.

\gesperrt{Integral Functions in an Algebraic Field.}
\end{center}

Already in the first volume we have had occasion several times to
introduce integral functions of an arbitrary number of independent
variables, and have also (in \S.\,141 f.)
%CONT>>

\vspace{1em}\hrule\vspace{0.5em}

\noindent 1878). Ueber die Discriminanten endlicher K\"orper (ebend.,
Bd.\,29, 1882). Ueber die Anzahl der Ideal-Classen in den verschiedenen
Ordnungen eines endlichen K\"orpers (Braunschweig 1877). Festschrift zur
S\"acularfeier des Geburtstages von \gesperrt{Gauss}. ``Zur Theorie der
Ideale'' and ``Ueber die Begr\"undung der Idealtheorie''. Nachrichten d.
Ges. d. Wissensch. in G\"ottingen 1894, 1895. To these belong also the
treatises of \gesperrt{Hilbert}, ``Ueber die Zerlegung der Ideale
etc.'', Mathem. Annalen, Bd.\,44, 1893. ``Grundz\"uge einer Theorie der
\gesperrt{Galois}'schen Zahlk\"orper'', G\"ottinger Nachrichten 1894.
\gesperrt{Hurwitz}, ``Zur Theorie der Ideale''. ``Ueber einen
Fundamentalsatz etc.'', G\"ottinger Nachrichten 1894, 1895.

\noindent\textsuperscript{1)} \gesperrt{Kronecker}, Grundz\"uge einer
arithmetischen Theorie der algebraischen Gr\"ossen. Festschrift zu
\gesperrt{Kummer's} 50j\"ahrigem Doctor-Jubil\"aum. Berlin 1882. (Also in
Bd.\,92 of \gesperrt{Crelle}'s Journal.) To be mentioned here are also the works of
\gesperrt{Hensel} in volumes 101, 103, 105, 111, 113 of \gesperrt{Crelle}'s Journal.

\noindent\textsuperscript{2)} Cf. the preface to the 4th edition of
\gesperrt{Dirichlet-Dedekind's} Vorlesungen \"uber Zahlentheorie.
